\documentclass[10pt,a4paper]{article}
\usepackage{fontspec}
\setmainfont{Liberation Sans}
\usepackage{polyglossia}
\setdefaultlanguage{polish}
\usepackage[a4paper,margin=19mm,headheight=15pt]{geometry}
\usepackage{xcolor,tabularx,booktabs,array,enumitem,fancyhdr,titlesec,hyperref}
\definecolor{navy}{HTML}{263650}
\definecolor{blue}{HTML}{5475A5}
\definecolor{pale}{HTML}{EFF3F8}
\definecolor{muted}{HTML}{5D6878}
\hypersetup{colorlinks=true,urlcolor=blue}
\pagestyle{fancy}
\fancyhf{}
\fancyhead[L]{\small\color{muted}Tutorland}
\fancyhead[R]{\small\color{muted}Stan prac}
\fancyfoot[C]{\small\color{muted}\thepage}
\renewcommand{\headrulewidth}{0.3pt}
\setlength{\parindent}{0pt}
\setlength{\parskip}{5pt}
\setlength{\emergencystretch}{3em}
\sloppy
\setlist[itemize]{leftmargin=*,itemsep=3pt,topsep=3pt}
\titleformat{\section}{\Large\bfseries\color{navy}}{\thesection.}{0.55em}{}
\newcolumntype{Y}{>{\raggedright\arraybackslash}X}
\begin{document}
\begin{titlepage}
\thispagestyle{empty}
\vspace*{32mm}
{\large\color{blue}\textbf{RAPORT POSTĘPU · 21 PAŹDZIERNIKA 2026}}
\vspace{12mm}
{\Huge\bfseries\color{navy}Tutorland\par}
\vspace{4mm}
{\LARGE Aktualny stan aplikacji\par}
\vspace{13mm}
\color{blue}\rule{\textwidth}{2pt}
\vspace{9mm}
{\large Przegląd wykonanej aplikacji, lokalnego API, plików danych i dokumentacji projektu.}
\vfill
\colorbox{pale}{\parbox{\dimexpr\textwidth-2\fboxsep\relax}{\textbf{W skrócie}\\
Działa katalog publiczny, profile korepetytorów, logowanie i rejestracja z wyborem roli, lokalne API z prywatnym zapisem CSV, rezerwacje, panele ucznia i korepetytora oraz obsługa opinii. Integracja Zoom OAuth i tworzenie spotkań 40-minutowych są zaimplementowane po stronie aplikacji, ale wymagają skonfigurowania aplikacji Zoom i sekretów środowiskowych.}}\par
\vspace{12mm}
{\small\color{muted}Raport opisuje kod projektu. Nie oznacza wdrożenia publicznego serwera ani aktywnego połączenia z zewnętrznym kontem Zoom.}
\end{titlepage}

\section{Zaimplementowane funkcje}
\begin{itemize}
\item \textbf{Katalog publiczny:} karty i profile korepetytorów, wyszukiwanie po imieniu i nazwisku, filtr przedmiotu i poziomu nauczania oraz podgląd dostępnych terminów przez gości.
\item \textbf{Konta i role:} osobne rejestracja i logowanie z wyborem ucznia lub korepetytora, sesja w ciasteczku \texttt{HttpOnly}, hasła przechowywane jako hashe scrypt oraz wylogowanie. Rola jest sprawdzana po stronie API.
\item \textbf{Lokalne API:} serwer Node obsługuje konta, katalog, profile, dostępność, rezerwacje i opinie. Zmienne dane trafiają do prywatnego katalogu \texttt{server/data/}; pliki CSV nie są serwowane statycznie. Dane demonstracyjne znajdują się w \texttt{server/seed/}.
\item \textbf{Rezerwacje:} uczeń może zapisać się na dostępny termin lekcji 40-minutowej. API weryfikuje rolę, przedmiot, dostępność, dni wolne, przeszłe godziny i nakładające się rezerwacje.
\item \textbf{Panel ucznia:} lista zaplanowanych i zakończonych spotkań, linki do spotkań Zoom, historia własnych opinii oraz formularze dodawania, edycji i usuwania. Opinię można dodać po zakończonej lekcji; obowiązuje jedna aktywna opinia na ucznia i korepetytora.
\item \textbf{Panel korepetytora:} edycja profilu, przedmiotów, poziomów i ceny; suwaki do ustawiania godzin dostępności w tygodniu; kalendarz liczby rezerwacji w bieżącym miesiącu i oznaczanie dni wolnych w kolejnym miesiącu.
\item \textbf{Powiadomienia korepetytora:} zaplanowane i zakończone lekcje oraz opinie uczniów. Korepetytor może je przeglądać, ale nie edytować.
\item \textbf{Połączenie Zoom:} przepływ OAuth korepetytora, szyfrowanie tokenów w prywatnym pliku, pobieranie zdjęcia profilu Zoom i tworzenie spotkań Zoom API trwających 40 minut. Uczeń otrzymuje link dołączenia, a link prowadzący jest dostępny tylko korepetytorowi.
\item \textbf{Dane konta demonstracyjnego:} usunięto wskazane konto testowe z aplikacji i źródłowych danych demonstracyjnych. Uczniowie nie mają zdjęć profilowych. Nazwę poziomu zmieniono na „Szkoła Podstawowa”.
\end{itemize}

\section{Konfiguracja i ograniczenia}
\small
\begin{tabularx}{\textwidth}{>{\bfseries}p{35mm}Y}
\toprule
Obszar & Stan i wymagane działanie\\
\midrule
Lokalny start & \texttt{npm run dev} uruchamia Vite i lokalne API. Zmienne dane pojawiają się w ignorowanym przez Git katalogu \texttt{server/data/}.\\
Zoom OAuth & Wymaga aplikacji OAuth w Zoom Marketplace, poprawnego adresu powrotu \texttt{/api/zoom/callback}, zakresu odczytu profilu i tworzenia spotkań oraz wartości z pliku \texttt{.env.example} zapisanych lokalnie w \texttt{.env}. Bez sekretów Zoom aplikacja jawnie informuje, że połączenie nie jest skonfigurowane.\\
Linki spotkań & Rezerwacja bez połączonego Zoom pozostaje w stanie oczekującym; po połączeniu konta serwer próbuje utworzyć oczekujące spotkania. Błędy Zoom są pokazywane w statusie rezerwacji.\\
Poczta e-mail & Aplikacja nie wysyła e-maili; powiadomienia są dostępne w panelach kont.\\
Środowisko & To lokalna aplikacja projektowa. Przed wdrożeniem publicznym potrzebna jest konfiguracja hostingu HTTPS, sekretów środowiskowych, kopii zapasowych danych i monitorowania.\\
\bottomrule
\end{tabularx}
\normalsize

\section{Sprawdzenia}
\begin{itemize}
\item Testy Node sprawdzają kodowanie CSV oraz podstawowe przepływy API: konta, sesje, dostęp do profilu, zapis rezerwacji, odrzucanie duplikatów i opinii przed odbytą lekcją.
\item \texttt{npm run lint} przechodzi bez ostrzeżeń.
\item \texttt{npm run build} generuje frontend produkcyjny.
\end{itemize}

\end{document}
